\section{Discussion}
\subsection{Interpretation}
The results highlight immense accuracy of traditional machine learning models, specifically Gradient Boosting. In particular, the machine learning model Gradient Boosting predict the students' final result with a balanced accuracy of just under $0.8$. This suggests that students' access to study materials, as well as study time, greatly affect the students' academic trajectory.

On the other hand, LLMs show a consistently poor performance, achieving a balanced accuracy of no further than $.5$. This is due to the fact that LLMs tend to witness a phenomenon of fallthrough to a default heuristics, most commonly ``Fail'' or ``Pass''. Interestingly, one-shot prompts perform better by a significant margin, but five-shot prompts perform no better than zero-shot, which may indicate overfitting.

From a pedagogical standpoint, our study serves as a rudimentary framework for implementating an early warning system for students. This cements the role of traditional machine learning models for such purpose and illustrates the unreliability of generative LLMs in classification problems.